% Generated draft: requires mathematical and source review.
\documentclass[11pt]{article}
\usepackage[margin=25mm]{geometry}
\usepackage{fontspec}
\setmainfont{Latin Modern Roman}
\usepackage{amsmath,amssymb,mathtools,unicode-math}
\setmathfont{Latin Modern Math}
\usepackage{enumitem,needspace,xurl,hyperref}
\setlength{\parindent}{0pt}
\setlength{\parskip}{5pt}
\setlength{\emergencystretch}{3em}
\newcommand{\N}{\mathbb{N}}
\newcommand{\Z}{\mathbb{Z}}
\newcommand{\Q}{\mathbb{Q}}
\newcommand{\R}{\mathbb{R}}
\newcommand{\C}{\mathbb{C}}
\newcommand{\F}{\mathbb{F}}
\newcommand{\A}{\mathbb{A}}
\newcommand{\PP}{\mathbb{P}}
\newcommand{\E}{\mathbb{E}}
\newcommand{\eps}{\varepsilon}
\newcommand{\dd}{\,\mathrm{d}}
\begin{document}
\section*{UnsolvedMath ID: 700374}
\textbf{Source problem code:} catalogue entry
\subsection{Definitions and mathematical statement}
Let \(G\) be a finite group and let \(A\leq G\) be an elementary abelian \(2\)-group, meaning a finite-dimensional vector space over \(\F_2\) under addition. Write \(A^g=g^{-1}Ag\) for \(g\in G\). Interpret the assertion that \(A\) is a trivial-intersection (\(TI\)) subgroup as
\[
\forall g\in G,\qquad A^g=A\quad\text{or}\quad A\cap A^g=\{1\}.
\]
Let \(P\) be a Sylow \(2\)-subgroup of \(G\) containing \(A\). Using the conventional interpretation of weak closure, put
\[
W_G(A,P)=\bigl\langle A^g\mid g\in G,\ A^g\leq P\bigr\rangle.
\]
The additional hypothesis is that, for some such \(P\), this subgroup is abelian:
\[
\forall x,y\in W_G(A,P),\qquad xy=yx.
\]
\textbf{Structural investigation requested by the source.} Describe the structure of the finite groups \(G\) admitting an elementary abelian \(TI\)-subgroup \(A\) and a Sylow \(2\)-subgroup \(P\geq A\) satisfying these conditions.

The source asks to investigate structure; it does not specify a conjectured conclusion, classification equivalence, or required form of description. The definitions of \(TI\) and weak closure above make their conventional meanings explicit, but are not supplied definitions. The record gives no rank bound for \(A\), no explicit nontriviality or properness requirement, and no hypothesis that conjugates of \(A\) generate \(G\). Whether additional conventions or hypotheses are intended requires review.
\subsection{Short English statement}
Investigate finite groups containing an elementary abelian \(2\)-subgroup whose distinct conjugates intersect trivially and whose weak closure in a containing Sylow \(2\)-subgroup is abelian. The source specifies a structural investigation without a proposed conclusion; its intended scope and subgroup conventions require mathematical review.
\subsection{Sources}
\begin{itemize}
\item Local UnsolvedMath registry, permanent ID \(700374\), exact source code \texttt{catalogue entry}, title \emph{Kourovka Notebook Problem 6.28}. Supplied upstream revision: \texttt{d371ee917ac9647fbb06cd6fc529deec6f3dda82}. The supplied registry labels the problem \emph{open}; this label has not been independently verified.
\item \emph{The Kourovka Notebook (Unsolved Problems in Group Theory)}, original problem 6.28, sixth issue (1978); supplied reference to the twenty-first issue edition (2026), \texttt{arXiv:1401.0300v48}. The supplied collection description names E. I. Khukhro and V. D. Mazurov as editors. This reference is supplied attribution and has not been independently consulted.
\item Proposed by A. A. Makhn\"ev, with proposed year 1978, according to the supplied record.
\end{itemize}
\end{document}
